\chapter{Banach空间}

\section{Banach空间的基本定义}


\begin{definition}[范数]
    设$X$是域$\mb{K}=\mb{R}$或$\mb{C}$上的线性空间，则$X$上的\bluebf{范数}是函数$\|\cdot\|\colon X\to[0,+\infty)$，满足：
    \begin{enumerate}[label=(\roman*)]
        \item $\|x\|=0\iff x=0$.
        \item $\|\lambda x\|=|\lambda|\cdot\|x\|$，$\forall\,\lambda\in\mb{K}$，$x\in X$.
        \item $\|x+y\|\leq\|x\|+\|y\|$，$\forall\, x,y\in X$.
    \end{enumerate}
\end{definition}
其中这三条性质分别称为范数的非退化性、一次齐次性和三角不等式.范数衡量了线性空间中元素的大小.\par
除了大小之外，为了研究收敛性、连续性等拓扑性质，我们需要为线性空间配备拓扑，其一的方式是为线性空间中的元素定义其距离.
\begin{definition}[度量]
    设$X$是一个拓扑空间，则$X$上的\bluebf{度量}是函数$\rho: X\times X\to[0,+\infty)$，满足：
    \begin{enumerate}[label=(\roman*)]
        \item $\rho(x,y)=0\iff x=y$.
        \item $\rho(x,y)=\rho(y,x)$，$\forall\, x,y\in X$.
        \item $\rho(x,z)\leq\rho(x,y)+\rho(y,z)$，$\forall\, x,y,z\in X$.
    \end{enumerate}
\end{definition}
度量就导出了一个拓扑空间中的距离的概念.\par
并且可以发现，由范数可以自然诱导出度量，因此任意赋范线性空间都是度量空间，其中$\rho(x,y):=\|x-y\|$.\par
反之则不一定成立，因为一个度量空间未必是线性的.例如所有的黎曼曲面都是度量空间，但黎曼曲线显然大部分都是非线性的.但我们可以研究由范数诱导出的度量空间都会满足什么样的性质.

\begin{remark}
    对于范数诱导出的度量$\rho$，它满足以下几个性质：
    \begin{itemize}
        \item 平移不变性：$\rho(x+w,y+w)=\rho(x,y)$
        \item 一次齐次性：$\rho(\lambda x,\lambda y)=|\lambda|\cdot\rho(x,y)$
    \end{itemize}\par
    前者让我们只需要考虑零点附近的性质就可以了解全局的性质，加上后者我们就只需要研究清除单位球就基本可以了解全局性质了，而关于单位球，有以下性质：
    \begin{itemize}
        \item 单位球是凸的：以后我们就记$\mb{B}_r(x)=\{y\in X:\rho(x,y)<r\},\mb{B}_r=\mb{B}_r(0),\mb{B}(x)=\mb{B}_1(x)$.\par
        \begin{sproof}
            给定$x,y\in\mb{B},\lambda\in(0,1)$，$\|\lambda x+(1-\lambda)y|\leq\|\lambda x\|+\|(1-\lambda)y\|=|\lambda|\|x\|+|1-\lambda|\|y\|<\lambda+1-\lambda=1$.\par
            $\Rightarrow \lambda x+(1-\lambda)y\in\mb{B}$.
        \end{sproof}
    \end{itemize}\par
    因此，当度量满足上述三个性质时，这个度量就和某个范数相容.\par
\end{remark}


赋予线性空间范数后，下面就可以研究上面的收敛性.\par

\begin{definition}[Cauchy列]
    设$X$是赋范线性空间，序列$\{x_n\}\subset X$是\bluebf{Cauchy列}$\iff \text{当}\, n,l\to\infty,\|x_n-x_l\|\to 0\iff \forall\,\varepsilon>0,\exists\, N\in\mb{N},\st\,\forall\, n,l>N,\|x_n-x_l\|<\varepsilon.$
\end{definition}

Cauchy列中的元素之间的距离会逐渐被控制得任意小，但一般的度量空间中Cauchy列未必收敛.例如在$(0,1)$中，序列$\{1-\frac{1}{n}\}$是Cauchy列，但它的极限点$1$不在$(0,1)$中.因此我们需要引入完备性的概念来保证Cauchy列的收敛性.

\begin{definition}[完备性]
    一个赋范线性空间$X$是\bluebf{完备的}当且仅当$X$中的每个Cauchy列都收敛.
\end{definition}

\begin{definition}[Banach空间]
    一个完备赋范线性空间称为\bluebf{Banach空间}.
\end{definition}

\section{Banach空间的例子}

\begin{instance}[$p-$范数]
    设$X=\left(\mb{R}^d,\|\cdot\|_p\right),1\leq p\leq \infty$.对$\vv{x}=\left(x^1,x^2,\cdots,x^d\right)\in\mb{R}^d$，定义其范数为$\|\vv{x}\|_p:=\begin{cases}
    \left(\sum_{i=1}^d|x^i|^p\right)^{\frac{1}{p}},&1\leq p<\infty\\
    \max_{1\leq i\leq d}|x^i|,&p=\infty
    \end{cases}$.\par
    下证上述确为范数.\par
    \begin{proof}
        (i)(ii)显然.\par
        (iii)当$p=\infty$时，显然得证.当$1\leq p<\infty$时，首先需要下列不等式.
        \begin{lemma}[Hölder不等式]
            $\left|\langle \vv{x},\vv{y}\rangle\right|=\left|\sum_{i=1}^d{x^iy^i}\right|\leq\|\vv{x}\|_p\cdot\|\vv{y}\|_q$，其中$1\leq p,q\leq\infty$满足$\frac{1}{p}+\frac{1}{q}=1$，称$p,q$\;\bluebf{Hölder共轭}.
        \end{lemma}
        \begin{sproof}
            若$\vv{x}=0$或$\vv{y}=0$，易证.\par
            否则不失一般性，设$\|\vv{x}\|_p=1=\|\vv{y}\|_q$.\par
            注意到，若$a,b\geq 0$，$\frac{1}{p}+\frac{1}{q}=1(1\leq p,q\leq\infty)$，则$ab\leq\frac{a^p}{p}+\frac{b^q}{q}$.\par
            因此$\left|\sum_{i=1}^d{x^iy^i}\right| \leq \sum_{i=1}^d\left(\frac{|x^i|^p}{p}+\frac{|y^i|^q}{q}\right) \xlongequal{\|\vv{x}\|_p=1=\|\vv{y}\|_q}\frac{1}{p}+\frac{1}{q}=1.$
        \end{sproof}
        下证三角不等式成立.\par
        $\setjot[-2]\begin{aligned}[t]
            \|\vv{x}+\vv{y}\|_p^p=\sum_{i=1}^d\|x^i+y^i\|^p &\leq \sum_{i=1}^d\left(|x^i|+|y^i|\right)^{p-1}\left(|x^i|+|y^i|\right) \\
            &= \sum_{i=1}^d\left(|x^i|+|y^i|\right)^{p-1}|x^i|+\sum_{i=1}^d\left(|x^i|+|y^i|\right)^{p-1}|y^i| \\
            &\leq \left(\sum_{i=1}^d\left(|x^i+y^i|^{p-1}\right)^q\right)^\frac{1}{q}\left(\left(\sum_{i=1}^d|x^i|^p\right)^\frac{1}{p}+\left(\sum_{i=1}^d|y^i|^p\right)^\frac{1}{p}\right) \\
            &=\left(\sum_{i=1}^d|x^i+y^i|^p\right)^\frac{p-1}{p}\left(\|\vv{x}\|_p+\|\vv{y}\|_p\right) \\
            &=\|\vv{x}+\vv{y}\|_p^{p-1}\left(\|\vv{x}\|_p+\|\vv{y}\|_p\right)
        \end{aligned}$\par
        移项得$\|\vv{x}+\vv{y}\|_p\leq \|\vv{x}\|_p+\|\vv{y}\|_p$.
    \end{proof}\par
    记$\mb{S}_X=\{x\in X:\|x\|=1\}$为单位球.当$d=2$时，我们可以画出不同$p$下的单位球的图像.\par
    \begin{figure}[H]
        \centering
        \vspace{-1 em}
        \includegraphics[width=0.4\textwidth]{fig/1.1.png}
        \captionof{figure}{不同$p$下的单位球}
    \end{figure}
    可以发现，当$p\geq 1$时，单位球都是凸集，并且若$p\neq 1$或$\infty$，单位球是一致凸的.因此处理$p-$范数时，有时会单独讨论$p=1$或$\infty$的情况.\par
    当$p<1$时，单位球不再是凸集，因为$p<1$时，对应定义的函数不再是范数.\par
\end{instance}

\begin{instance}[$\ell^p$空间]
    设$1\leq p\leq \infty$，可以定义数列组成的空间$\ell^p$为$\{\{\alpha_j\}_{j\in\mb{N}}:\|\{\alpha_j\}\|_p<\infty\}=\left(\mb{R}^\infty,\|\cdot\|\right)$，并定义其范数为$\|x\|_p:=\begin{cases}
        \left(\sum_{j=1}^\infty|\alpha_j|^p\right)^{\frac{1}{p}},&1\leq p<\infty\\
        \sup_{j\in\mb{N}}|\alpha_j|,&p=\infty
    \end{cases}$.\par
    可以证明上述空间为一个赋范线性空间.\par
\end{instance}

\begin{instance}[连续函数空间]
    设$K$为紧致Hausdorff空间，定义$C(K)=\{f:K\to\mb{R}:f\text{连续}\}$，并定义其范数为上极限范数$\|f\|=\sup_{x\in K}|f(x)|<\infty$.范数的有限性来自于紧集上的连续函数是有界的，因此上述空间为一个赋范线性空间.\par
    下面证明它是完备的.
    \begin{proof}
        取$\{f_n\}\subset C(K)$为Cauchy列，则当$n,\ell\to\infty$时，$\|f_n-f_\ell\|=\sup_{x\in K}\left|f_n(x)-f_\ell(x)\right|\to 0$.\par
        对任意固定的$x\in K$，$\{f_n(x)\}$是$\mb{R}$中的Cauchy列.\par
        由于$\mb{R}$是完备的，因此$\exists\, f:K\to\mb{R},f_n(x)\to f(x)$.\par
        由于上述极限为一致极限，$f_n$都连续，因此$f$也连续.\par
        下证$f_n$一致收敛到$f$，则其在$C(K)$中收敛.\par
        已知当$n,\ell\to\infty$时，$\|f_n-f_\ell\|=\sup_{x\in K}\left|f_n(x)-f_\ell(x)\right|\to 0$.\par
        令$\ell\to\infty$，得当$n\to\infty,\left|f_n(x)-f(x)\right|\to 0$，进而当$n\to\infty,\|f_n-f\|\to 0$.\par
        因此$\{f_n\}$在$C(K)$中收敛到$f$，$C(K)$是完备的.
    \end{proof}
    现在可以研究$C(K)$中单位球的性质.\par
    令$K=[0,1],f_n=x^n$.则$\|f_n\|=\sup_{x\in[0,1]}|x^n|=1$，因此$\{f_n\}\subset\mb{S}_{C(K)}$.\par
    但我们知道$f_n\to\chi_{1}$，而$\chi_{1}$不是连续函数，因此$\{f_n\}$在$C(K)$中没有收敛子列.\par
    从而$C(K)$中的单位球不是紧集.\par
\end{instance}
从上面一个例子中可以看出，连续函数空间的单位球是一个有界闭集，但它不是紧集.但在实数域等有限维线性空间中，有界闭集都是紧集.因此我们可以得到以下结论：
\begin{theorem}
    设$X$是一个赋范线性空间，则有界闭集是紧集$\iff \dim X<\infty$.
\end{theorem}
这个定理预示着有限维线性空间的拓扑性质未必能推广到无限维线性空间之中，下一个例子同样展示了这点.\par
\begin{instance}
    设$X=C([0,1])$，$Y=\left\{[0,1]\,\text{上多项式}\right\}$，则$\overline{Y}=X\supsetneq Y$，因此$Y$不是闭集
\end{instance}
因此，对于无限维赋范空间而言，子空间未必是闭集，我们需要额外的条件.
\begin{theorem}
    设$(X,\|\cdot\|)$是一个Banach空间，设$Y\leq X$是一个线性子空间，则$(T,\|\cdot\|)$是Banach空间$\iff$ $Y$是闭集.
\end{theorem}
\begin{proof}
    $\Rightarrow$：设$(Y,\|\cdot\|)$是Banach空间，取任意$\{y_n\}\subset Y,\st y_n\to x\in X$.\par
    由于$\{y_n\}$收敛，则$\{y_n\}$是$Y$中的Cauchy列，因此$\{y_n\}\to y\in Y$.\par
    由于极限的唯一性，$x=y\in Y$，因此$Y$是闭集.\par
    $\Leftarrow$:设$Y$是闭集，取$\{y_n\}\subset Y$为一个Cauchy列.\par
    因为$(X,\|\cdot\|)$是Banach空间，$y_n\to x\in X$.\par
    又因为$\{y_n\}\subset Y$, $Y$是闭集，因此$x\in Y$，从而$(Y,\|\cdot\|)$是Banach空间.
\end{proof}

下面继续详细研究数列空间$\ell^p$.\par
\begin{instance}
    设$X=\ell^2$，$e_i=(0,\cdots,0,1,0,\cdots)$，其中$1$在第$i$个坐标.\par
    $Z:=\mathrm{span}\,\{e_i:i\in\mb{N}\}=\left\{\sum_{i=1}^N\alpha_i e_i:N\in\mb{N},\alpha_i\in\mb{R}\right\}$.\par
    考虑$x=\left(1,\frac{1}{2},\frac{1}{3},\frac{1}{4},\cdots\right)$，则$x\notin Z$.取$x_n=\left(1,\frac{1}{2},\cdots,\frac{1}{n},0,\cdots,0,\cdots\right)$，则$x_n\to x$，就有$x\in\overline{Z}$.\par
    并且$\|x_n-x\|_{\ell^2}=\left\|\left(0,\cdots,0,\frac{1}{n+1},\frac{1}{n+2},\cdots\right)\right\|_{\ell^2}=\left(\sum_{j=n+1}^\infty\frac{1}{j^2}\right)^{\frac{1}{2}}\xlongrightarrow{n\to\infty}0$.\par
    因此$\left(Z,\|\cdot\|_{\ell^2}\right)$不是Banach空间.\par
    事实上，$Z=\left\{(\alpha_j):\exists\, N\in\mb{N},\st\,\alpha_j=0,\forall\,j\geq N\right\}$且$\overline{Z}=\ell^2$，这意味着$Z$是$\ell^2$中的一个稠密子集.\par
    更进一步，可以证明$Z$上的任意范数均不完备.
\end{instance}

\begin{definition}[级数收敛]
    $\sum_n{x_n}$在$X$中收敛当且仅当$S_n=\sum_{n=1}^N{x_n}$收敛.\par
    $\sum_n{x_n}$在$X$中绝对收敛当且仅当$\sum_n\|x_n\|$收敛.\par
\end{definition}
下面可以证明以下结论.
\begin{theorem}[]
   设$(X,\|\cdot\|)$是赋范线性空间，则$(X,\|\cdot\|)$是Banach空间$\iff$ $X$中每个绝对收敛序列都收敛. 
\end{theorem}
\begin{proof}
    $\Rightarrow$：设$X$是Banach空间，设$\sum\|x_n\|<\infty$，考虑$S_n=\sum_{i=1}^n{x_i}$.\par
    $\|S_n-S_\ell\|=\|\sum_{j=\ell+1}^n{x_j}\|\leq\sum_{j=\ell+1}^n\|x_j\|\xlongrightarrow{n,\ell\to\infty} 0$.\par
    因此$\{S_n\}$是Cauchy列，故收敛，因而$\sum_n{x_n}$收敛.\par
    $\Leftarrow$：任取Cauchy列$\{y_n\}\subset X$，可取一个迅速收敛的子列$\{y_{n_k}\}\st\,\|y_{n_{k+1}}-y_{n_k}\|<2^{-k}$.\par
    那么$\sum_k\left(y_{n_{k+1}}-y_{n_k}\right)$绝对收敛，因此由假设收敛.\par
    故$y_{n_k}=y_{n_1}+\sum_{j=1}^{k-1}\left(y_{n_{j+1}}-y_{n_j}\right)$收敛.\par
    \begin{claim*}
        任意有收敛子列的Cauchy列也收敛.
    \end{claim*}
\end{proof}
最终可以证明上面所定义的数列空间是完备的.
\begin{theorem}[Riesz-Fischer]
    $\ell^p(\mb{R}),\ell^p(\mb{C})$是Banach空间，$\forall\, 1\leq p\leq \infty$.
\end{theorem}
\begin{proof}
    \begin{enumerate}[label=(\roman*)]
        \item \textbf{\boldmath$p=\infty$}\par
        取$\{x_n\}$是$\ell^\infty$中Cauchy列，其中$x_n=\{x_n^{(i)}\}_{i\in\mb{N}}$.\par
        则当$n,k\to\infty$时，$\|x_n-x_k\|_{\ell^\infty}=\sup_{i\in\mb{N}}\left|x_n^{(i)}-x_k^{(i)}\right|\to 0$.\par
        对每个固定的$i\in\mb{N}$，$\{x_n^{(i)}\}_{n\in\mb{N}}$是$\mb{K}=\mb{R}$或$\mb{C}$下的Cauchy列，因而收敛.\par
        因此当$n\to\infty$，$x_n^{(i)}\to y^{(i)}$.\par
        \begin{claim}{$y=\{y^{(i)}\}$是$\ell^\infty$范数下$\{x_n\}$的极限}
            $\setjot[-2]\begin{aligned}[t]
                \|x_n-y\|_{\ell^\infty} &=\sup_{i\in\mb{N}}\left|x_n^{(i)}-y^{(i)}\right|=\sup_{i\in\mb{N}}\left|x_n^{(i)}-\lim_{k\to\infty}x_k^{(i)}\right| \\
                &=\sup_{i\in\mb{N}}\lim_{k\to\infty}\left|x_n^{(i)}-x_k^{(i)}\right|\to 0,\;\text{当}\, n\to\infty
            \end{aligned}$\par
        \end{claim}
        又可知$y\in \ell^\infty$，故$\ell^\infty$是Banach空间.
        \item \textbf{\boldmath$1\leq p<\infty$}\par
        设$\{x_n\}\subset \ell^p$是Cauchy列，其中$x_n=\{x_n^{(i)}\}_{i\in\mb{N}}$.\par
        则$\forall\,\varepsilon>0,\exists\, N\in\mb{N},\st\forall\, n,k\geq N,\|x_n-x_k\|_{\ell^p}^p=\sum_{i=1}^\infty\left|x_n^{(i)}-x_k^{(i)}\right|^p<\varepsilon$.\par
        和$\ell^\infty$的情形相同，当$n\to\infty$时，有$x_n^{(i)}\to y^{(i)}$.\par
        那么对任意$M\in\mb{N}$，有$\sum_{i=1}^M \left|x_n^{(i)}-y^{(i)}\right|^p=\lim_{k\to\infty}\sum_{i=1}^M\left|x_k^{(i)}-y^{(i)}\right|^p\leq\varepsilon$.\par
        因此对于足够大的$n$，$\sum_{i=1}^\infty \left|x_n^{(i)}-y^{(i)}\right|\leq\varepsilon$，此处$\varepsilon>0$是任意的.\par
        所以$x_n\to y=\{y^{(i)}\}_{i\in\mb{N}}$在$\ell^p$中收敛.
    \end{enumerate}
\end{proof}


\section{线性算子}

赋范线性空间的另一研究方向是赋范线性空间之间的线性算子，其中更关注于有界线性算子.\par
\begin{definition}[有界线性算子]
    设$X,Y$是赋范线性空间，线性算子$T:X\to Y$是\bluebf{有界的}当且仅当它满足以下任意条件：
    \begin{enumerate}[label=(\roman*)]
        \item $\exists\, M>0,\st\forall\, x\in X.\|Tx\|_Y\leq M\|x\|_X$.
        \item $\|T\tilde{x}\|_Y\leq M,\forall\,\tilde{x}\in\mb{S}_X$
        \item $\|T\tilde{x}\|_Y\l  eq M,\forall\,\tilde{x}\in\mb{B}^X$
        \item $T\left(\overline{\mb{B}_1^X(0)}\right)\subseteq \overline{\mb{B}_M^Y(0)}$
    \end{enumerate}
\end{definition}
上述四个条件虽然是彼此之间等价的，但它们分别从不同的角度描述了有界线性算子.第一个条件是从纯粹分析的角度描述，它给出了范数的估计.而到第四个条件则是从拓扑的角度出发，描述了单位球在映射下的像所满足的条件.由于泛函分析有强烈的拓扑背景，因此在分析和拓扑的视角中相互切换是学习泛函分析的重要技巧.\par
\begin{definition}[算子代数]
    设$X,Y$是赋范线性空间，线性算子$T:X\to Y$的\bluebf{算子代数}定义为
    \[\|T\|_{X\to Y}=\|T\|:=\inf\{M>0:\|Tx\|_Y\leq M\|x\|_X,\forall\, x\in X\}\]
    $T$是有界线性算子当且仅当$\|T\|_{X\to Y}<\infty$.
\end{definition}
我们记$\mcal{L}(X,Y):=\{T:X\to Y\,\text{线性算子}\}$为$X\to Y$的所有的线性算子，$\mcal{B}(X,Y):=\{T\in\mcal{L}(X,Y):\|T\|<\infty\}$为$X\to Y$的所有的有界线性算子.可以证明$\left(\mcal{B}(X,Y),\|\cdot\|_{X\to Y}\right)$是一个赋范线性空间.自然地，我们想知道$\mcal{B}(X,Y)$何时是一个Banach空间.

\begin{definition}[]
    设$X,Y$是赋范线性空间，则$\mcal{B}(X,Y)$是Banach空间$\iff$ $Y$是Banach空间.
\end{definition}
\begin{proof}
    设$\{T_n\}\subset \mcal{B}(X,Y)$是Cauchy列，则当$n,\ell\to\infty$时，$\|T_n-T_\ell\|\to 0$.\par
    对$\forall\, x\in X$，$\|T_n x-T_\ell x\leq \|T_n-T_\ell\|\|x\|$，所以$\{T_n x\}$是$Y$中的Cauchy列.\par
    因此$T_n x\to Tx$，其中$T:X\to Y$是算子.\par
    有$T_n$的线性性可知，$T$是线性算子，下证$T\in\mcal{B}(X,Y)$.\par
    由于$\{T_n\}$为Cauchy列，$\exists\, N\in\mb{N},\st\,\forall\, n\geq N,\|T_n-T_{N_0}\|\leq 1$.\par
    因此$\setjot[-2]\begin{aligned}[t]
        \|Tx\| &\leq \|T_{N_0}x\|+\lim_{n\to\infty}\|T_n x-T_{N_0} x\| \\
        &\leq \|T_{N_0}x\|+1\cdot\|x\| \\
        &\leq \left(\|T_{N_0}\|+1\right)\|x\|,\forall\, x\in X
    \end{aligned}$\par
    即$\|T\|\leq \|T_{N_0}\|+1<\infty$.\par
    $\|T_n-T\|=\sup\left\{\|T_n x-Tx\|:x\in X,\|x\|\leq 1\right\}=\sup_{\|x\|\leq 1}\lim_{k\to\infty}\|T_n x-T_k x\|\xlongequal{n\to \infty}0$
\end{proof}

可以观察到，若$\{T_n\}\subset\mcal{B}(X,Y)$，$T_n$在范数下收敛到$T$，则$\|T_n\|\to\|T\|$以及$\forall\, x\in X,T_n x\to Tx$.\par
但对于后者，其逆命题并不成立，即$\forall\, x\in X,T_n x\to Tx$并不意味着$T_n$在范数下收敛到$T$.例如设$X=Y=\ell^2,e_i=(0,\cdots,0,1,0,\cdots)$，其中$1$在第$i$个坐标，定义$T_n e_i=\begin{cases}
    e_i,&i\neq n\\
    2e_n,&i=n
\end{cases}$，则$\forall\, x\in X,T_n x\to x$，但$\|T_n-I\|=\sup_{\|x\|\leq 1}\|T_n x-x\|\geq \|T_n e_n-e_n\|=\|e_n\|=1$，因此$T_n$在范数下不收敛到$I$.这表明前者的条件比后者更强，我们有时也称前者为\bluebf{强范数收敛}.\par
下面一个定理表明在线性算子的前提下，有界和连续是等价的.\par
\begin{theorem}[]
    设$X,Y$是赋范线性空间，$Y\in\mcal{L}(X,Y)$，则$T$连续$\iff$$T$有界.
\end{theorem}
\begin{proof}
    \begin{claim*}
        $T$连续当且仅当$T$在$x=0$处连续
    \end{claim*}
    $\Rightarrow$：设$T$在$x=0$处连续，则$\forall\,\varepsilon>0,\exists\, \delta>0,\st\,\forall\,\|x\|<\delta,\|Tx\|<\varepsilon$.\par
    令$\tilde{x}=\frac{x}{\delta}$，那么$\forall\, \|\tilde{x}\|\geq 1,\|T\tilde{x}\|\leq\frac{\varepsilon}{\delta}$\par
    $\Rightarrow \|T\|\leq \frac{\varepsilon}{\delta}$.\par
    $\Leftarrow$：设$T$有界，则$\forall\, \|x\|\leq 1,\|Tx\|\leq M$.\par
    则$\forall\, \varepsilon>0$，令$\delta=\frac{\varepsilon}{M}$，那么$\forall\,\|x\|<\delta,\|Tx\|\leq\varepsilon$，故$T$连续.
\end{proof}\par

容易想到，有界线性算子之间的复合仍旧是有界线性算子.
\begin{definition}[]
    设$X,Y,Z$是赋范线性空间，若$T\in\mcal{B}(X,Y),S\in\mcal{B}(Y,Z)$，则$ST\equiv S\circ T\in\mcal{B}(X,Z)$.
\end{definition}
\begin{proof}
    $\setjot[-2]\begin{aligned}[t]
        \|ST\|_{X\to Z} &=\sup\left\{\|S(Tx)\|_Z:\forall\, x\in X,\|x\|\leq 1\right\} \\
        &\leq \sup\left\{\|S\|_{Y\to Z}\|Tx\|_Y:\forall\, x\in X,\|x\|\leq 1\right\} \\
        &=\|S\|_{Y\to Z}\sup\left\{\|Tx\|_Y:\forall\, x\in X,\|x\|\leq 1\right\} \\
        &=\|S\|_{Y\to Z}\|T\|_{X\to Y}
    \end{aligned}$\par
    即$\|ST\|\leq \|S\|\|T\|$.
\end{proof}
特别地，如果$X=Y$，则记$\mcal{B}(X,X)=\mcal{B}(X)$，并且可以定义其中的乘法$\setjot[-6]\begin{aligned}[t]
    \mcal{B}(X)\times\mcal{B}(X) &\to\mcal{B}(X) \\
    (T,S) &\mapsto S\circ T
\end{aligned}$，并且满足$\|ST\|\leq \|S\|\|T\|$.如果$X$为一个Banach空间，则$\mcal{B}(X)$也是一个Banach空间且在上方有乘法，称为\bluebf{Banach代数}.\par
\begin{instance}
    $n$阶矩阵是一个Banach代数.\par
    $\left(C[0,1],\|\cdot\|_\mathrm{sup}\right)$，$(fg)(x)=f(x)g(x)$是一个Banach代数.
\end{instance}
\begin{instance}[有界线性算子]
    \begin{enumerate}[label=(\arabic*)]
        \item $A\in Mat(n\times k;\mathbb{R})\cong \mcal{L}(\mathbb{R}^k,\mathbb{R}^n)$
        \item $\setjot[-6]\begin{aligned}[t]
            \Lambda:\ell^p &\to\ell^p \\
            (x_1,x_2,\cdots) &\mapsto (\lambda_1 x_1,\lambda_2 x_2,\cdots)
        \end{aligned} $\par
        \begin{claim*}
            $\Lambda$是有界线性算子$\iff$ $\{\lambda_j\}\in\ell^\infty$.
        \end{claim*}
        并且，如果考虑$\setjot[-6]\begin{aligned}[t]
            \ell^\infty &\xrightarrow{\Phi} \mcal{B}(\ell^p) \\ 
            \{\lambda_j\}&\mapsto \Lambda
        \end{aligned}$，那么$\Phi$是等距同构，即$\|\Phi\left(\{\lambda_j\}\right)\|_{\mcal{B}(\ell^\infty)}=\|\{\lambda_j\}\|_{\ell^\infty}$.
        \begin{proof}
            $"\leq "$：$\setjot[-2]\begin{aligned}[t]
                \|\Phi\left(\{\lambda_j\}\right)\|_{\mcal{B}(\ell^p)}&=\sup\left\{\left(\sum_{j=1}^\infty|\lambda_jx_j|^p\right)^\frac{1}{p}:\left(\sum_j^\infty|x_j|^p\right)^\frac{1}{p}\leq 1\right\}\\
                &\leq\sup\left\{(\sup_j|\lambda_j|)\left(\sum_j^\infty|x_j|^p\right)^\frac{1}{p}:\left(\sum_j^\infty|x_j|^p\right)^\frac{1}{p}\leq 1\right\} \\
                &=\|\{\lambda_j\}\|_{\ell^\infty}
            \end{aligned}$ \par
            $"\geq "$：要证$\forall\,\varepsilon>0,\|{\lambda_j}\|_{\ell^\infty}\leq \|\Phi\left({\lambda_j}\right)\|_{\mcal{B}(\ell^p)}+\varepsilon$.\par
            取$j_0\in\mb{N}$满足$|\lambda_{j_0}|\geq \|\{\lambda_j\}\|_{\ell^\infty}-\varepsilon$.\par
            则$\setjot[-2]\begin{aligned}[t]
                \|\Phi\left(\{\lambda_j\}\right)\|_{\mcal{B}(\ell^p)}&=\sup\left\{\|\Lambda x\|_{\ell^p}:x\in\ell^p,\|x\|_{\ell^p}\leq 1\right\} \\
                &\geq \|\Lambda e_{j_0}\|_{\ell^p},e_{j_0}=(0,\ldots,1,\ldots,0) \\
                &=|\lambda_{j_0}|\geq \|\{\lambda_j\}\|_{\ell^\infty}-\varepsilon
            \end{aligned}$\par
        \end{proof}
        \item 设$X=\ell^p,1\leq p\leq \infty$，定义左右位移算子分别为$L:X\to X,L(x_1,x_2,x_3,\cdots)=(x_2,x_3,\cdots),R:X\to X,R(x_1,x_2,x_3,\cdots)=(0,x_1,x_2,\cdots)$.\par
        则$\|L\|=1=\|R\|$，并且$LR=\mathrm{Id}_X\neq RL$，这表明有界线性算子可能并不存在逆算子，而要区分左逆和右逆.\par
        本例中的原因是$R$是单射，$R$是满射，因此$RL$不单不满，必然不是恒等映射.\par
        但是我们可以发现$\dim (X/\ran R)=1,\dim\ker L=1$，被称为\bluebf{Fredholm算子}，它是泛函分析中可逆算子的推广.
        \item 设$\setjot[-2]\begin{aligned}[t]
            I:C[0,1]&\to C[0,1] \\ 
            f&\mapsto F(x)=\int_0^x{f(y)\d y}
        \end{aligned}$\par
        则$\forall\, x\in[0,1],\left|\int_0^xf(y)\d y\right|\leq\int_0^x|f(y)|\d y\leq \left(\int_0^x 1\d y\right)\|f\|_{C[0,1]}=x\|f\|_{C[0,1]}$.\par
        故$\|F\|_{C[0,1]}\leq \|f\|_{C[0,1]}$，从而$I\in\mcal{B}\left(C[0,1]\right)$的算子范数为$\|I\|=1$(在$f=1$时取等).
    \end{enumerate}
    下面再看几个数列空间的例子，记$C:=\left\{{\alpha_j}\in\mb{C}\,\text{或}\,\mb{C}:\,\text{当}\,j\to\infty,\alpha_j\,\text{收敛}\right\}$.$C_0:=\left\{{\alpha_j}\in C:\lim_{j\to\infty}\alpha_j=0\right\}$，它是$\|\cdot\|_{\ell^\infty}$下的Banach空间.$C_{00}:=\left\{{\alpha_j}\in C:\,\text{有限项}\,\alpha_j\,\text{非零}\right\}$.
    \begin{enumerate}[label=(\arabic*),resume]
        \item 设$X=\left(C_{00},\|\cdot\|_{\ell^\infty}\right)$.\par
        考虑$\setjot[-2]\begin{aligned}[t]
            T:X&\to X \\
            (a_1,a_2,a_3,\cdots) &\mapsto \left(a_1,\frac{a_2}{2},\frac{a_3}{3},\cdots\right)
        \end{aligned}$\par
        由于$T$收缩，$T\in\mcal{B}(X)$且$\|T\|=1$.\par
        $T$明显是一个双射，但$T^{-1}\in\mcal{L}(X)$，且$T^{-1}(b_1,b_2,b_3,\cdots)=(b_1,2b_2,3b_3,\cdots)$，因此$\|T^-1\|=\infty$，故$T^{-1}\notin\mcal{B}(X)$.
    \end{enumerate}
\end{instance}
上述的最后两例表明，有界线性算子不一定存在逆算子，并且即使其存在逆算子，其逆算子也未必是一个可逆线性算子，关于这一点，有以下定理阐述.
\begin{theorem}[]
    设$X$是Banach空间，$T\in\mcal{B}(X)$是双射，则$T^{-1}\in\mcal{B}(X)$.
\end{theorem}

从上述定理中也可以看出，$(C_{00},\|\cdot\|)$不是一个Banach空间.\par
接下来更深入的研究有界线性算子的可逆性.在线性空间中，我们定义过特征值的概念，若$\lambda$为特征值，则$\exists\, v\neq 0,\st Tv=\lambda v\iff \ker(\lambda I-T)\neq\{0\}\Rightarrow \lambda I-T$可逆.注意在有限维线性空间中，最后一个箭头是双向的，逆命题成立，但在无限维线性空间中，逆命题不成立.因此我们需要一个新的概念来描述有界线性算子的可逆性.
\begin{definition}[谱]
    设$X$是Banach空间，$T\in\mcal{B}(X)$，$T$的\bluebf{谱}是$\sigma(T)=\Spec(T):=\{\lambda\in\mb{C}:\lambda I-T\,\text{不可逆}\}$.
\end{definition}
谱的一个重要性质是
\begin{theorem}[]
    $\sigma(T)$是$\mb{C}$中的非空紧集.
\end{theorem}
这个定理的证明需要后续学习的工具，我们目前可以证明$\sigma(T)$一定是一个闭集.\par
定义$\Aut(X):=\{T\in\mcal{B}(X):T^{-1}\,\text{存在}\}$.回到有限维线性空间的情形举例.对$\Aut(\mb{R}^d)=GL(\mb{R}^d)$.由于行列式是连续的，因此行列式等于零的集合是一个闭集，从而$\Aut(\mb{R}^d)$是一个开集.对于一般的Banach空间，我们也可以证明$\Aut(X)$是$\mcal{B}(X)$中的开集，从而$\sigma(T)$是$\mb{C}$中的闭集.\par
\begin{theorem}[]
    $\Aut(X)$是$\mcal{B}(X)$的开集.
\end{theorem}
\begin{proof}
    \begin{enumerate}[label=Step \arabic*.,leftmargin=4em]
        \item $\mb{B}_1(\mathrm{Id}_X)\subset \Aut(X)$，即$\forall\, T\in\mcal{B}(X)$满足$\|T\|< 1$，$I-T$可逆.\par
        \begin{claim}{$(I-T)^{-1}=\sum_{k=0}^{\infty}T^k$}
            $\sum_{k=0}^\infty\|T^k\|\leq \sum_{k=0}^\infty \|T\|^k<\infty$.\par
            所以$\sum T^k$在$\mcal{B}(X)$中绝对收敛，故$\sum T^k$收敛.\par
            令$S=\sum_{k=0}^\infty T^k\in\mcal{B}(X)$，则当$n\to\infty$，$S_n=\sum_{k=0}^nT^k\to S$.\par
            由于$S_n$只有有限项，因此$S_n(I-T)=(I-T)S_n=T^0-T^{n+1}\xlongrightarrow{n\to\infty}I$.\par
            再令$n\to\infty$，有$S(I-T)=(I-T)S=I$，因此$S=(I-T)^{-1}$.
        \end{claim}
        \item 取任意$T\in\Aut(X)$.\par
        当$\|T^{-1}S\|<1$时，$T-S=T(I-T^{-1}S)$可逆，即$\|S\|<\|T^{-1}\|^{-1}$.\par
        则$\exists\, \delta=\|T^{-1}\|^{-1},\st\forall\, S\in\mcal{B}(X),\|S\|<\delta,T-S\in \Aut(X)$
    \end{enumerate}\par
    这说明$\forall\, T\in\Aut(X),\exists\,\mb{B}_\delta(T)\subset \Aut(X)$，从而$\Aut(X)$是开集.
\end{proof}
类似的想法，定义预解集的概念.
\begin{definition}[预解集]
    设$X$是Banach空间，$T\in\mcal{B}(X)$，则$\rho(T):=\{\lambda\in\mb{C}:\lambda I-T\in\Aut(X)\}$称为$T$的\bluebf{预解集}.
\end{definition}
由定理1.8，可以证明预解集也是一个开集，从而$\sigma(T)$是$\mb{C}$中的闭集\par
\begin{corollary}[]
    $\rho(T)$是$\mcal{B}(X)$中的开集，$\sigma(T)$是$\mb{C}$中的闭集.
\end{corollary}
\begin{proof}
    定义$\setjot[-6]\begin{aligned}[t]
        f:\mb{C} &\to \mcal{B}(X) \\
        \lambda &\mapsto \lambda I-T
    \end{aligned}$，显然$f$是一个连续映射.\par
    由于$\Aut(X)$是$\mcal{B}(X)$中的开集，所以$\rho(T)=f^{-1}(\Aut(X))$是$\mb{C}$中的开集.\par
    可以看出$\sigma(T)$是$\rho(T)$的补集，因此取补得到$\sigma(T)$是$\mb{C}$中的闭集.
\end{proof}

在这一节的最后，引入泛函的概念.
\begin{definition}[有界线性泛函]
    设$X$是Banach空间，记$X^*=\mcal{B}(X,\mb{R})$或$\mcal{B}(X,\mb{C})$
\end{definition}





